%chapter introduce un nuevo capítulo
\chapter{Resumen}

Las enfermedades neurodegenerativas representan un desafío significativo en el campo de la medicina, ya que su progresión y los efectos en los pacientes son altamente complejos y variados. Desde el Alzheimer hasta el Parkinson, estas enfermedades afectan no solo la calidad de vida de los pacientes, sino también su capacidad para llevar una vida independiente. Uno de los mayores retos en el tratamiento de estas enfermedades es la predicción del declive cognitivo, que puede variar ampliamente entre individuos y dificulta el desarrollo de estrategias terapéuticas efectivas.\\

En este contexto, el presente estudio se centra en el desarrollo de modelos predictivos de progresión de la enfermedad de Parkinson utilizando el entorno Leaspy en Python. A partir de una cohorte de 423 pacientes obtenidos de la Parkinson's Progression Markers Initiative (PPMI), se busca abordar este desafío crucial al predecir con precisión la aparición y evolución del declive cognitivo en pacientes con diferentes estados de la enfermedad.\\

Se emplea una estrategia de entrenamiento de modelos utilizando una parte de la cohorte, mientras que la otra se reserva para evaluar la precisión de las predicciones. El objetivo primordial es identificar la combinación óptima de biomarcadores que permita anticipar con exactitud el avance cognitivo de los pacientes. Para lograr esto, se implementan modelos progresivos para explorar todas las posibles combinaciones de variables relevantes, aprovechando la paralelización a través de plataformas como Curie y Magerit para agilizar el proceso de análisis y modelado.\\

Al comparar y evaluar los modelos generados, se analizan diversas métricas de desempeño de los DPMs, incluyendo la precisión en la estimación temporal y la capacidad de diagnóstico a lo largo del tiempo, así como la predicción de la conversión hacia el declive cognitivo. Este enfoque integrado no solo proporciona una comprensión más profunda de la progresión de la enfermedad de Parkinson, sino que también ofrece perspectivas clave que pueden contribuir significativamente al desarrollo de nuevos enfoques terapéuticos y estrategias de tratamiento en el futuro.

\paragraph{Palabras clave:} Leaspy, DPM, PPMI, biomarcadores, paralelización, Curie, Magerit, declive cognitivo, Parkinson.

\chapter{Abstract}

Neurodegenerative diseases pose a significant challenge in the field of medicine, as their progression and effects on patients are highly complex and varied. From Alzheimer's to Parkinson's, these diseases affect not only patients' quality of life but also their ability to live independently. One of the biggest challenges in treating these diseases is predicting cognitive decline, which can vary widely among individuals and complicates the development of effective therapeutic strategies.\\

In this context, the present study focuses on developing predictive models of Parkinson's disease progression using the Leaspy environment in Python. Drawing from a cohort of 423 patients obtained from the Parkinson's Progression Markers Initiative (PPMI), the aim is to address this crucial challenge by accurately predicting the onset and evolution of cognitive decline in patients at different stages of the disease.\\

A model training strategy is employed using one part of the cohort, while the other part is reserved for evaluating prediction accuracy. The primary objective is to identify the optimal combination of biomarkers that can precisely anticipate patients' cognitive advancement. To achieve this, progressive models are implemented to explore all possible combinations of relevant variables, leveraging parallelization through platforms like Curie and Magerit to streamline the analysis and modeling process.\\

Upon comparing and evaluating the generated models, various performance metrics of the DPMs are analyzed, including accuracy in temporal estimation and diagnostic capability over time, as well as prediction of conversion to cognitive decline. This integrated approach not only provides a deeper understanding of Parkinson's disease progression but also offers key insights that can significantly contribute to the development of new therapeutic approaches and treatment strategies in the future.

\paragraph{Keywords:} Leaspy, PPMI, biomarkers, parallelization, Curie, Magerit, cognitive decline, Parkinson.